\section*{Chapter \ref{ch:b1:continuity} --- Limits and Continuity}

\begin{solution}{exo:b1:continuity:1}
Take $u_n = \frac{1}{2\pi n + \pi/2}$ and $v_n = \frac{1}{2\pi n}$:
both tend to $0^+$, yet $\sin\frac{1}{u_n} = 1$ and
$\sin\frac{1}{v_n} = 0$. Two sequences, two different limits of
images: by \cref{thm:b1:continuity:seqchar}, no limit at $0^+$.

$x \sin\frac1x$: squeezed by $\abs{x\sin\frac1x} \leq \abs x \to 0$,
so the limit at $0$ exists and equals $0$.
\end{solution}

\begin{solution}{exo:b1:continuity:2}
$f = \lfloor\cdot\rfloor$ is continuous on $\R \setminus \Z$
(locally constant) and discontinuous at each $n \in \Z$: left limit
$n - 1$, value $n$.

$g(x) = x - \lfloor x\rfloor$: same discontinuity points (the
identity is continuous, so $g$ inherits the jumps of $f$); at $n \in
\Z$, left limit $1 \neq 0 = g(n)$.

$h(x) = \lfloor x\rfloor + (x - \lfloor x\rfloor)^2$: on $\intco{n}{n
+ 1}$, $h(x) = n + (x - n)^2$, continuous there; at $x = n$, the left
limit is $(n-1) + 1 = n = h(n)$: the jumps cancel. $h$ is continuous
on $\R$ (and strictly increasing).
\end{solution}

\begin{solution}{exo:b1:continuity:3}
$P(x) = x^5 - 3x + 1$: $P(-2) = -32 + 6 + 1 = -25 < 0$; $P(0) = 1 >
0$; $P(1) = -1 < 0$; $P(2) = 27 > 0$. Three sign changes on the
disjoint segments $\intcc{-2}{0}$, $\intcc{0}{1}$, $\intcc{1}{2}$:
by \cref{thm:b1:continuity:ivt}, at least three roots. (Being of
degree $5$, $P$ has at most five; a variation study would show
exactly three.)
\end{solution}

\begin{solution}{exo:b1:continuity:4}
Let $P = a_{2m+1}X^{2m+1} + \dots$ with $a_{2m+1} > 0$ (else replace
$P$ by $-P$). Factoring the dominant term, $P(x) = a_{2m+1}
x^{2m+1}\bigl(1 + o(1)\bigr)$ as $x \to \pm\infty$: so $P(x) \to
+\infty$ at $+\infty$ and $-\infty$ at $-\infty$. Pick $a$ with
$P(a) < 0$ and $b$ with $P(b) > 0$: the intermediate value theorem on
$\intcc{a}{b}$ provides a root.
\end{solution}

\begin{solution}{exo:b1:continuity:5}
Let $g(x) = f(x) - x$, continuous on $\intcc{0}{1}$. Since $f$ maps
into $\intcc{0}{1}$: $g(0) = f(0) \geq 0$ and $g(1) = f(1) - 1 \leq
0$. By \cref{thm:b1:continuity:ivt}, $g(c) = 0$ for some $c$: a
fixed point.

Necessity of the hypotheses: on $\intcc{0}{1}$, the discontinuous map
$f(x) = 1$ for $x \leq \frac12$, $f(x) = 0$ for $x > \frac12$ has no
fixed point; on the interval $\intoo{0}{1}$ (not a segment), $f(x) =
\frac x2$ is continuous into $\intoo{0}{1}$ with no fixed point
(the candidate $0$ is missing); on $\R$, $f(x) = x + 1$.
\end{solution}

\begin{solution}{exo:b1:continuity:6}
$g(x) = \cos x - x$ is continuous, $g(0) = 1 > 0$, $g(1) = \cos 1 -
1 < 0$: a solution exists in $\intoo{0}{1}$
(\cref{thm:b1:continuity:ivt}). Uniqueness: $g$ is strictly
decreasing on $\R$ --- for $x \leq 0$, $g(x) \geq 1 - x > 0$ has no
zero anyway; and $g'(x) = -\sin x - 1 \leq 0$ with equality only at
isolated points ($x \equiv -\frac\pi2 \bmod 2\pi$), so $g$ is
strictly decreasing (\cref{ch:b1:derivative}; alternatively: on
$\intcc{0}{1}$, $\cos$ is strictly decreasing and $-x$ too, so $g$
is). A strictly monotone function vanishes at most once.
\end{solution}

\begin{solution}{exo:b1:continuity:7}
Fix $M = f(0) + 1$. There is $A > 0$ with $f(x) \geq M$ for $\abs x
\geq A$ (definition of the two infinite limits; take the larger
threshold). On the segment $\intcc{-A}{A}$, the extreme value theorem
(\cref{thm:b1:continuity:evt}) provides $c$ with $f(c) =
\inf_{\intcc{-A}{A}} f \leq f(0)$. For $\abs x \geq A$: $f(x) \geq M
> f(0) \geq f(c)$. So $f(c)$ is the global minimum.
\end{solution}

\begin{solution}{exo:b1:continuity:8}
On the segment $\intcc{0}{T}$, $f$ is bounded and attains its bounds
(\cref{thm:b1:continuity:evt}); by periodicity, these are the bounds
on all of $\R$, still attained.

Let $g(x) = f(x + \frac T2) - f(x)$, continuous. Then
\[
g(0) + g\bigl(\tfrac T2\bigr)
= \bigl(f(\tfrac T2) - f(0)\bigr) + \bigl(f(T) - f(\tfrac T2)\bigr)
= f(T) - f(0) = 0 :
\]
$g(0)$ and $g(\frac T2)$ have opposite signs (or one vanishes), so
the intermediate value theorem on $\intcc{0}{T/2}$ gives $c$ with
$g(c) = 0$, i.e.\ $f(c + \frac T2) = f(c)$.
\end{solution}

\begin{solution}{exo:b1:continuity:9}
First the inequality: for $0 \leq y \leq x$,
\[
\bigl(\sqrt y + \sqrt{x - y}\bigr)^2 = x + 2\sqrt{y(x-y)} \geq x,
\]
so $\sqrt x \leq \sqrt y + \sqrt{x - y}$, i.e.\ $\sqrt x - \sqrt y
\leq \sqrt{x - y}$. Hence $\abs{\sqrt x - \sqrt y} \leq
\sqrt{\abs{x - y}}$ for all $x, y \geq 0$.

Uniform continuity: given $\varepsilon > 0$, take $\delta =
\varepsilon^2$; then $\abs{x - y} \leq \delta$ implies $\abs{\sqrt x
- \sqrt y} \leq \sqrt\delta = \varepsilon$.

Not Lipschitz near $0$: $\frac{\sqrt x - \sqrt 0}{x - 0} =
\frac{1}{\sqrt x} \to +\infty$ as $x \to 0^+$, so no constant $k$
can dominate all difference quotients.
\end{solution}

\begin{solution}{exo:b1:continuity:10}
Suppose $f$ is injective, continuous, and not strictly monotonic.
Then there are $a < b < c$ in $I$ with $f(b)$ not between $f(a)$ and
$f(c)$ --- indeed, if for \emph{all} triples the middle value were
between the outer ones, $f$ would be monotonic (compare any two
pairs; a short case check). Say $f(b) > \max(f(a), f(c))$ (the other
case is symmetric, replace $f$ by $-f$). Choose $v$ with $\max(f(a),
f(c)) < v < f(b)$. By the intermediate value theorem applied on
$\intcc{a}{b}$ and on $\intcc{b}{c}$, there are $u_1 \in
\intoo{a}{b}$ and $u_2 \in \intoo{b}{c}$ with $f(u_1) = v = f(u_2)$:
two distinct points with equal images, contradicting injectivity.
\end{solution}

\begin{solution}{exo:b1:continuity:11}
$f(0) = f(0+0) = 2f(0)$ gives $f(0) = 0$; $f(-x) = -f(x)$ from $0 =
f(x - x)$. Set $\alpha = f(1)$. Induction: $f(n) = n\alpha$ for $n
\in \N$, then for $n \in \Z$ by oddness. For $q \in \N^*$:
$q\,f(\frac pq) = f(p) = p\alpha$ (add $\frac pq$ to itself $q$
times), so $f(\frac pq) = \alpha\frac pq$: $f = \alpha\,
\mathrm{id}$ on $\Q$.

Let now $x \in \R$ and $(r_n)$ a sequence of rationals with $r_n \to
x$ (density, \cref{thm:b1:reals:density}, applied in nested
intervals; or $r_n = \frac{\lfloor nx\rfloor}{n}$). Continuity:
$f(x) = \lim f(r_n) = \lim \alpha r_n = \alpha x$.
\end{solution}

\begin{solution}{exo:b1:continuity:12}
Let $\varepsilon > 0$. By the limit at $+\infty$, there is $A$ with
$\abs{f(x) - \ell} \leq \frac\varepsilon2$ for $x \geq A$; hence for
$x, y \geq A$: $\abs{f(x) - f(y)} \leq \varepsilon$ (no closeness
needed).

On the segment $\intcc{0}{A + 1}$, Heine's theorem
(\cref{thm:b1:continuity:heine}) gives $\delta_0 > 0$ for this
$\varepsilon$; set $\delta = \min(\delta_0, 1)$.

Now take any $x, y \geq 0$ with $\abs{x - y} \leq \delta$, say $x
\leq y$. If $y \leq A + 1$: both lie in the segment, and the Heine
$\delta_0$ applies. Otherwise $y > A + 1$, and then $x \geq y - 1 >
A$: both lie in $\intco{A}{+\infty}$, where the limit argument
applies. In both cases $\abs{f(x) - f(y)} \leq \varepsilon$: uniform
continuity.
\end{solution}
